\section{Discussion and Future Work}
\vspace{-0.5em}
In this work, we propose two different verification methods, based on robust scenario optimization and conformal prediction, to provide probabilistic safety guarantees for neural reachable tubes.
Our methods allow a direct trade-off between resilience to outlier errors in the neural tube, which are inevitable in a learning-based approach, and the strength of the probabilistic safety guarantee.
Furthermore, we show that split conformal prediction, a widely used method in the machine learning community for uncertainty quantification, reduces to a scenario-based approach, making the two methods equivalent not only for verification of neural reachable tubes but also more generally.
We hope that our proof will lead to future insights into the close relationship between the highly related but disparate fields of conformal prediction and scenario optimization.
Finally, we propose an outlier-adjusted verification approach that harnesses information about the error distribution in neural reachable tubes to recover greater safe volumes.
We demonstrate the efficacy of the proposed approaches for the high-dimensional problems of multi-vehicle collision avoidance and rocket landing with no-go zones.
Altogether, these are important steps toward using learning-based reachability methods to compute safety assurances for high-dimensional systems in the real world.

In the future, we will explore how the key idea of the outlier-adjusted verification approach, using cost labels as a supervised learning signal, can be used to enhance the accuracy of learning-based reachability methods like DeepReach.
Other directions include providing safety assurances in the presence of worst-case disturbances and in real-time for tubes that are generated online.